\section{Limitations}\label{sec:limitations}

\paragraph{Thermal limit of one card.} The four cards sit in one host with a shared air path, and the card at PCIe
bus 21 runs hottest. With the engine bitstream loaded, that card heats above 84\,\textdegree C within minutes even
when no kernel runs. Every measurement therefore runs under a guard: before each repetition, a card at 85\,\textdegree C
(FPGA) or 80\,\textdegree C (HBM) is parked on the small bitstream of the baseline design and cooled below
84\,\textdegree C; a run during which a card reaches 90\,\textdegree C (FPGA) or 85\,\textdegree C (HBM) is kept,
flagged, and ends that configuration. Table~\ref{tab:thermal} lists the \NumThermalStops~configurations that the guard
stopped, all on card 21; where a stopped configuration was run again later, the complete run is the one reported. In one early attempt the cooling wait was done
with the engine bitstream still loaded; card 21 reached 96\,\textdegree C (FPGA) and 91\,\textdegree C (HBM), its AXI
firewall tripped and the driver reset the card. That run is excluded from all results and the procedure was changed
as described. The widest runs, at $n=31$ on four cards, succeeded only after cooling pauses; on this host the time
available for long runs at full width is limited by the temperature of one card, not by memory.

\paragraph{Routing limit of the engine count.} A build with sixteen complex128 engines did not route
(Section~\ref{sec:architecture}). The engine uses eight compute units with four pseudo channels each; it reaches
all 32 pseudo channels of a card, but the peak kernel throughput is that of eight engines.

\paragraph{Half of the HBM holds the state.} At full width the state occupies 8~GiB of the 16~GiB of each card
(Section~\ref{subsec:res-maxwidth}): every pseudo channel holds $2^{24}$ amplitudes (256~MiB) of a slice next to the
engine's instruction segment, which shares the buffer, and the next power of two would fill the 512~MiB channel
completely. Storing the instructions in a separate buffer, on a further HBM port or in on chip memory, would remove this
limit and give one more qubit per configuration ($n=30$ on one card, $n=32$ on four); we did not build this variant.

\paragraph{Data movement across cards goes through the host.} Gates that need amplitudes of another card are
executed after an exchange through host memory over PCIe. Peer to peer transfers between the cards are disabled on
this host (the platform reports that they would need 32~GB of PCIe address space per card), and we did not change that system
setting. The cost of this path for each circuit family is given in
Sections~\ref{subsec:res-scaling} and~\ref{subsec:res-ablation}. This path is why more cards do not make a run faster:
on four cards the exchanges take \NumScaleExchQftFour\,\% of the end to end time of the QFT at $n=28$ and
\NumScaleExchRandFour\,\% for the random circuit, and they set the time of every circuit whose slice qubits carry non
diagonal gates. Peer to peer transfers or a faster interconnect would shorten exactly this part.

\paragraph{Card count.} The partition assigns the top $\log_2 C$ qubits to the card index, so the number of cards
must be a power of two. The runtime accepts any such count, but only $C = 1$, 2 and 4 were run, because the host has
four cards. Eight cards would need a host with eight x16 slots, and every further card qubit adds remaps across cards
for the circuits that act on it; with the exchange through host memory, the trend of Section~\ref{subsec:res-scaling},
a growing exchange share with every doubling of the cards, would be expected to continue. We make no measured claim for
eight or more cards.

\paragraph{One host, one platform.} All measurements come from one host with four U55C cards. The behaviour of the
exchange on hosts with a different PCIe topology or with peer to peer transfers enabled was not measured.
